\documentclass[11pt, oneside]{article}   	% use "amsart" instead of "article" for AMSLaTeX format
\usepackage{geometry}                		% See geometry.pdf to learn the layout options. There are lots.
\geometry{letterpaper}                   		% ... or a4paper or a5paper or ... 
%\geometry{landscape}                		% Activate for rotated page geometry
%\usepackage[parfill]{parskip}    		% Activate to begin paragraphs with an empty line rather than an indent
\usepackage{graphicx}				% Use pdf, png, jpg, or eps§ with pdflatex; use eps in DVI mode
								% TeX will automatically convert eps --> pdf in pdflatex		
\usepackage{amssymb}

%SetFonts

%SetFonts
\usepackage{hyperref}
\usepackage{xcolor}

\title{Housing Allocation}
\author{Nick Arnosti}
%\date{}							% Activate to display a given date or no date

\begin{document}
\maketitle
%\section{}
%\subsection{}



Post a). Inefficiency! And a solution.
-Most housing agencies in US punish rejections, leading to inefficient matching. This inefficiency is especially acute when there is a large imbalance.
-Not punishing rejections would lead to better matching




Post b). Why do they do that???

Answers:
-Maybe haven?t thought about it carefully. I hope this is the case.
-Administrative overhead from many declined offers (+delay). Solution: Amsterdam (have them apply)
-Don?t want people to have choice (forced integration). This isn?t guaranteed to work (may not move), and is mean ? why can?t poor people have nice things?
-Want to discourage those that don?t really need it. This is somewhat legitimate ? if many people have very specific desires, this might prevent them from matching (relative to those with more flexibility). However, there are caveats. First, you could end up with a situation where everyone accepts everything (i.e. your housing is nice, or people who don?t need it aren?t applying). Second, inflexibility might not mean low need ? maybe just means very constrained! Third, if you are really worried about all these people that don?t need housing, find another way to identify them (i.e. income).

Post c). What about other systems?
-New York uses a lottery
-Equivalence
-CEEI/virtual currency
-Common lottery (relate to school choice work)


\end{document}  